% Ad Atticum 7.26 --- headnote
% ad-atticum-07-26, 15 February 49 BC, Formian villa

\textit{Cicero to Atticus, written from the Formian
villa on the fifteenth day before the Kalends of March
49 BC (the manuscript dateline: \textit{Scr. in
Formiano xv K. Mart. a. 705 (49)}). This is the
closing letter of \textit{Ad Atticum} book 7; with book
8 the picture darkens again. For a moment here it has
brightened: news has come in from Rome of Domitius's
stand and of the Picene cohorts; Afranius too is well
reported; Caesar's threats are being met with quotation
rather than panic.}

\textit{Section 1 catches that brief lift --- the
flight being prepared has been checked; Caesar's
interdict, an Ennian-tasting line about being found
``on a second day's light,'' is being thrown back at
him. Section 2 is Cicero answering Atticus on his own
political standing. The two-edged sentence ---
\textit{ego me ducem in civili bello quoad de pace
ageretur negavi esse, non quin rectum esset sed quia
quod multo rectius fuit id mihi fraudem tulit} --- is
the self-knowledge of a man who has chosen one position
over another and is paying for it. The quoted phrase
\textit{pro tuis rebus gestis amplissimis} is the
formula in which Pompey, before all this, would have
moved Cicero's second consulship and triumph.}

\textit{Section 3 closes book 7 with the running
business --- a money matter Terentia has answered;
Dionysius the tutor and his obligations; the boys
(Marcus and the younger Quintus) likely to winter at
Formiae --- and then the personal decision, made
plain: \textit{si enim erit bellum, cum Pompeio esse
constitui}. ``If there is to be a war, I have decided
to be with Pompey.'' What follows that line, in book
8, is the trying-out of the decision.}
